%% ma: L10-B02-0012
%% lop: 10
%% chuong: 1
%% bai: 2
%% dang: 10.2.7
%% loai: TN4
%% muc: VD
%% dap_an: "B"
%% nguon: "Ảnh giáo viên gửi 10/10/2026 (ThS. Nguyễn Hoàng Việt, Chương 2 Tập hợp), A-Đề 1, Phần 1 Câu 10"
%% kien_thuc: "Bài 2 (tập hợp và các phép toán trên tập hợp)"
%% trang_thai: cho-duyet
%% ly_do: "Dạng toán mới đề xuất 10.2.7: tìm tham số để các tập hợp thỏa quan hệ (giao khác rỗng, tập con, rời nhau)"
\begin{ex}
Cho tập $M=\{x\in\mathbb R\mid(x+3)(x-1)\ge0\}$ và $N=(m-1;m+1]$. Tìm $m$ để $M\cap N\ne\varnothing$.
\choice{$\left[\begin{array}{l}m\le-2\\ m>0\end{array}\right.$}{\True $\left[\begin{array}{l}m<-2\\ m\ge0\end{array}\right.$}{$\left[\begin{array}{l}m\le-2\\ m\ge0\end{array}\right.$}{$-2\le m<0$}
\loigiai{$M=(-\infty;-3]\cup[1;+\infty)$. $M\cap N=\varnothing$ khi và chỉ khi $N\subset(-3;1)$, tức $m-1\ge-3$ và $m+1<1$, hay $-2\le m<0$. Vậy $M\cap N\ne\varnothing\Leftrightarrow m<-2$ hoặc $m\ge0$.}
\end{ex}
